\section{能力 Vignettes：产品把多个模型能力组合成可完成的工作}

接下来五组画面不是功能清单，而是在展示 \term{compositionality（可组合性）}。一个有用产品通常同时调用语言理解、代码生成、执行环境、视觉渲染、记忆和迭代反馈。单独测某个能力的 benchmark 很高，仍可能在组合链路中失败；因此每个 vignette 都应被读成一个端到端任务图。

\subsection{从生成代码到个人工具}

Stripe 风格 dashboard 例子把自然语言需求变成前端代码和可交互页面。这里至少存在四个验收层：需求约束是否被提取、代码是否能运行、视觉层级是否可读、用户修改后是否保持既有行为。任何一层失败，产品都不能只用“代码 token 正确率”解释。

\lecturefigure{slide-06-stripe-dashboard.jpg}{自然语言生成 Stripe 风格 dashboard：代码与界面同时成为输出}{00:04:41--00:05:00}

\lecturefigure{slide-07-correlation-explorer.jpg}{面向非数学用户的交互式相关性探索器}{00:05:00--00:05:51}

相关性探索器进一步要求系统把概念转成可操纵对象。若用户拖动数据分布，却看不到相关系数、散点结构与反例同步变化，交互只是在“演示动画”。好的产品 eval 应记录：核心变量是否可操作、反馈是否即时、例外情况是否可观察、用户是否能从操作中修正原有误解。

\begin{importantbox}{端到端任务成功不是子指标相加}
设需求解析、代码正确、执行成功、视觉可读和用户完成任务的事件分别为 $A_1,\dots,A_5$。端到端成功是联合事件
\[
P(\text{success})=P(A_1\cap A_2\cap A_3\cap A_4\cap A_5),
\]
而不是五个准确率的平均。若某一层是硬门槛，即使其他层接近满分，整体仍会被最弱环节限制。
\end{importantbox}

\subsection{图像、移动端与零成本软件}

手绘草图到生成图像的画面强调另一种 interface contract：用户提供低保真意图，模型补全风格与细节。输出漂亮不代表遵循了构图、对象关系和作者意图；评测需要把 aesthetic quality、instruction adherence、identity preservation 与 editability 分开。

\lecturefigure{slide-08-sketch-to-image.jpg}{从手绘草图到生成图像：模型补全细节，但作者仍需控制方向}{00:05:51--00:06:02}

\lecturefigure{slide-09-generative-mobile-game.jpg}{移动端即时生成小游戏：内容、代码与触控交互被组合}{00:05:52--00:06:24}

\lecturefigure{slide-10-zero-cost-software.jpg}{“软件创建趋近零成本”的 demo：移动界面由对话即时生成}{00:06:24--00:07:14}

这两张移动端画面支持的是“原型边际成本下降”，不是生产软件已经免费。真实交付仍包含需求澄清、测试、数据安全、维护、分发与责任。尤其当模型生成可执行软件时，必须把执行环境隔离：代码生成能力越强，越不能把任意网络、文件和凭证权限默认交给模型。

\teachervoice{00:04:04--00:07:36，Karina 把个人 dashboard、游戏、图像和移动软件串成“增强创造力”的证据，并希望未来系统更主动、更个性化。讲义保留她的愿景，同时补上工程边界：主动性必须和权限、可撤销性、成本与验证共同设计。}

\begin{warningbox}{“零成本软件”最容易漏算什么}
生成一次 UI 的 token 成本可能很低，但 total cost of ownership 还包括需求错误、测试、运行基础设施、第三方依赖、权限审计、长期维护和事故恢复。产品研究应测完成真实任务的总成本，而不是只报一次生成价格。
\end{warningbox}

\subsection{本章小结}

能力 vignette 的共同结构是：用户意图经过模型与工具链，被转成可观察、可编辑、可运行的产物。它们证明了新的 product surface 有价值，也暴露了 benchmark 之外的失败面。下一章开始把这些失败面写成两个 scaling paradigm 和两条产品研究路径。

